\subsection{Anomaly cancellation $\Rightarrow$ counterterm restoration in EG (strong-closure)}
\label{app:eg_anomaly_counterterm_restoration}

\AuditTag This appendix turns anomaly cancellation from a static filter (Appendix~\ref{app:anomaly_theorem_filters}) into an explicit strong-closure statement about the existence of counterterms restoring BRST/Slavnov--Taylor identities in the EG carrier.

\subsubsection{Finite renormalization freedom as local counterterms}
\paragraph{Finite renormalization group.}
In EG, different choices of extensions on the diagonals correspond to finite renormalizations.
At fixed perturbative order and fixed EFT truncation, these freedoms are parameterized by a finite-dimensional space of local functionals compatible with power counting and symmetries.

\paragraph{Counterterm action.}
Let $R$ denote a finite renormalization transformation acting on time-ordered products (equivalently on the S-matrix functional).
Applying $R$ shifts the interaction by local counterterms within the admissible class.

\subsubsection{BRST cohomology and obstruction class}
\paragraph{Consistency condition.}
If the ST identity fails by $\Delta$ (Appendix~\ref{app:eg_brst_st_consistency}), the nilpotency of the BRST differential implies a consistency relation for $\Delta$ (Wess--Zumino type), placing $\Delta$ into a BRST cohomology class at a fixed ghost number.

\paragraph{Removal criterion.}
If the relevant cohomology class is trivial in the admissible local functional space, then there exists a local counterterm that cancels $\Delta$ and restores ST to the stated order.
If nontrivial, the anomaly is an unavoidable obstruction.

\subsubsection{Strong-closure theorem statement (template)}
\begin{theorem}[Anomaly cancellation implies ST-restoring counterterms (template)]
\label{thm:eg_anomaly_implies_restoration_template}
Fix a truncation order $N$ (operator dimension and/or loop order) and an EG perturbative order.
Assume the anomaly filters in Appendix~\ref{app:anomaly_theorem_filters} hold for the declared chiral content and gauge group.
Then any local breaking term $\Delta$ of the ST identity compatible with the truncation is BRST-exact within the admissible local functional space, hence removable by a finite renormalization; consequently, the ST identity can be restored to the stated order.
\end{theorem}

\AuditTag Theorem~\ref{thm:eg_anomaly_implies_restoration_template} is a strong-closure \emph{template}: it becomes a completed proof once the admissible local functional space is fixed (by power counting and field content) and the cohomology computation is reduced to the standard anomaly coefficients recorded in Appendix~\ref{app:anomaly_theorem_filters}.

\subsubsection{Interface to the EFT representative and error budget}
\InterfaceTag Restoring ST by local counterterms is the theorem-level mechanism that justifies treating Ward/BRST identities as exact \emph{to the declared truncation order}.
Any residual mismatch beyond that order is attributed to the truncation remainder and is budgeted in the EFT error decomposition (Appendix~\ref{app:eft_error_bounds}).

